\documentclass[11pt,a4paper]{article}
\usepackage[T1]{fontenc}
\usepackage[utf8]{inputenc}
\usepackage{amsmath,amssymb,xcolor}
\usepackage[margin=25mm]{geometry}
\usepackage[colorlinks=true,urlcolor=blue]{hyperref}
\definecolor{nexorablue}{HTML}{2459B5}
\setlength{\parindent}{0pt}
\setlength{\parskip}{7pt}
\setlength{\emergencystretch}{2em}
\renewcommand{\familydefault}{\sfdefault}
\begin{document}

{\small\bfseries\color{nexorablue} NEXORA CPGE}\par
{\small MP / KUT-SPE-07-01 / Corrigé}\par
{\LARGE\bfseries Diagonalisation avec une valeur propre double}\par
\vspace{0.25em}\hrule\vspace{0.6em}

\section*{Spectre}
Le développement suivant la première ligne donne $(X-2)((X-4)^2-4)=(X-2)^2(X-6)$.

\section*{Calcul des noyaux}
Pour $\lambda=2$, les équations se réduisent à $x+2y-z=0$. Les vecteurs $(-2,1,0)$ et $(1,0,1)$ forment une base de ce plan. Pour $\lambda=6$, on trouve $x=0$ et $z=-2y$, donc $(0,1,-2)$ engendre la droite propre.

\section*{Changement de base}
Les trois vecteurs forment une base puisque le plan propre et la droite propre sont en somme directe. On vérifie la relation $AP=PD$ sur les trois colonnes.
\[ P=\begin{pmatrix}-2&1&0\\1&0&1\\0&1&-2\end{pmatrix},\qquad D=\operatorname{diag}(2,2,6). \]

\bigskip
{\footnotesize Adaptation française par Nexora CPGE. Source : Ken Kuttler (Lyryx),
\href{https://math.libretexts.org/Bookshelves/Linear_Algebra/A_First_Course_in_Linear_Algebra_(Kuttler)/07\%253A_Spectral_Theory/7.02\%253A_Diagonalization}{\textit{A First Course in Linear Algebra}}, § 7.2, exemple 1.
\href{https://creativecommons.org/licenses/by/4.0/}{CC BY 4.0}.}
\end{document}
